\documentclass[10pt]{article}

% =========================================================
% TEMPLATE
% =========================================================

\input{~/Dropbox/Latex/TemplateExam/template-exam.tex}
% Si el template está en la misma carpeta del examen, usar:
% % =========================================================
% TEMPLATE MAESTRO PARA EXÁMENES
% template-exam.tex
% =========================================================

% =========================================================
% PAQUETES
% =========================================================

\usepackage{fontspec}
\usepackage{geometry}
\usepackage{xcolor}
\usepackage{graphicx}
\usepackage{amsmath}
\usepackage{amssymb}
\usepackage{mathtools}
\usepackage{array}
\usepackage{booktabs}
\usepackage{tabularx}
\usepackage{multicol}
\usepackage[inline]{enumitem}
\usepackage{changepage}
\usepackage{tikz}
\usepackage{etoolbox}
\usepackage{lastpage}

\usetikzlibrary{
    arrows.meta,
    positioning,
    calc,
    shapes.geometric
}

% =========================================================
% COLORES
% =========================================================

\definecolor{ExamGreen}{HTML}{1B6B45}
\definecolor{ExamGreenLight}{HTML}{EAF4EE}
\definecolor{ExamGray}{HTML}{5F6368}
\definecolor{ExamLightGray}{HTML}{F4F5F5}
\definecolor{ExamLine}{HTML}{CED4D1}
\definecolor{ExamBlack}{HTML}{202124}

% Alias para reutilizar diagramas de otras clases
\providecolor{BrandBlue}{HTML}{2864A5}
\providecolor{BrandGreen}{HTML}{1B6B45}
\providecolor{BrandYellow}{HTML}{D9A514}

% =========================================================
% INFORMACIÓN INSTITUCIONAL
% =========================================================

\providecommand{\CourseName}{Nombre de la asignatura}
\providecommand{\CourseCode}{Código}
\providecommand{\Faculty}{Facultad}
\providecommand{\Institution}{Institución}
\providecommand{\Instructor}{Nombre del docente}
\providecommand{\Email}{correo@institucion.edu.co}

\providecommand{\ExamTitle}{Examen}
\providecommand{\ExamTopics}{Temas evaluados}
\providecommand{\AcademicPeriod}{Periodo académico}
\providecommand{\ExamDuration}{60 minutos}
\providecommand{\ExamValue}{100 puntos}
\providecommand{\ExamQuestions}{20}

% Dejar únicamente el nombre del archivo.
% Ejemplo:
% \renewcommand{\InstitutionLogo}{LogoPoli.png}
\providecommand{\InstitutionLogo}{}

% =========================================================
% CONFIGURACIÓN GENERAL
% =========================================================

\setlength{\parindent}{0pt}
\setlength{\parskip}{0pt}

\setlist[itemize]{
    leftmargin=5mm,
    itemsep=0.5mm,
    topsep=1mm,
    parsep=0pt,
    partopsep=0pt
}

\setlist[enumerate]{
    itemsep=0.5mm,
    topsep=1mm,
    parsep=0pt,
    partopsep=0pt
}

% =========================================================
% IMÁGENES
% =========================================================

% LaTeX buscará imágenes en estas carpetas
\graphicspath{
    {./}
    {Template/}
    {Images/}
}

% ---------------------------------------------------------
% MINI LOGO
% ---------------------------------------------------------
% Uso:
%
% \renewcommand{\InstitutionLogo}{LogoPoli.png}
%
% El archivo puede estar en:
% ./LogoPoli.png
% Template/LogoPoli.png
% Images/LogoPoli.png
% ---------------------------------------------------------

% ---------------------------------------------------------
% MINI LOGO
% ---------------------------------------------------------

\newcommand{\MiniLogo}{
    \ifx\InstitutionLogo\empty
    \fbox{
        \parbox[c][0.9cm][c]{1.8cm}{
            \centering
            \scriptsize
            \color{ExamGray}
            Sin logo
        }
    }
    \else
    \IfFileExists{\InstitutionLogo}{
        \includegraphics[
        height=1.5cm,
        keepaspectratio
        ]{\InstitutionLogo}
    }{
        \fbox{
            \parbox[c][0.9cm][c]{1.8cm}{
                \centering
                \scriptsize
                \color{ExamGray}
                Logo no\\
                encontrado
            }
        }
    }
    \fi
}

% ---------------------------------------------------------
% IMAGEN PARA PREGUNTAS
% ---------------------------------------------------------
% Uso:
%
% \ExamImage{robot.png}
%
% o
%
% \ExamImage[4cm]{robot.png}
% ---------------------------------------------------------

% ---------------------------------------------------------
% IMAGEN PARA PREGUNTAS
% ---------------------------------------------------------
% Uso:
%
% \ExamImage{robot.png}
%
% o
%
% \ExamImage[4cm]{robot.png}
% ---------------------------------------------------------

\newcommand{\ExamImage}[2][5cm]{
    \begin{center}
        \IfFileExists{#2}{
            \includegraphics[
            width=#1,
            keepaspectratio
            ]{#2}
        }{
            \fbox{
                \parbox[c][1.5cm][c]{0.75\linewidth}{
                    \centering
                    \scriptsize
                    \color{ExamGray}
                    Imagen no encontrada\\[1mm]
                    \texttt{\detokenize{#2}}
                }
            }
        }
    \end{center}
}

% =========================================================
% ESPACIADOS CONFIGURABLES
% =========================================================

\newlength{\QuestionBefore}
\newlength{\QuestionAfter}
\newlength{\AnswerItemSep}
\newlength{\AnswerTopSep}
\newlength{\AnswerIndent}
\newlength{\InlineAnswerSep}
\newlength{\GridAnswerSep}

% Valores por defecto
\setlength{\QuestionBefore}{1mm}
\setlength{\QuestionAfter}{0.7mm}
\setlength{\AnswerItemSep}{0.5mm}
\setlength{\AnswerTopSep}{0.8mm}
\setlength{\AnswerIndent}{0.45cm}
\setlength{\InlineAnswerSep}{0.7cm}
\setlength{\GridAnswerSep}{2mm}

% =========================================================
% MODO COMPACTO
% =========================================================

\newcommand{\ExamCompact}{
    \geometry{
        letterpaper,
        left=1.15cm,
        right=1.15cm,
        top=0.8cm,
        bottom=0.8cm
    }

    \setlength{\parindent}{0pt}
    \setlength{\parskip}{0pt}
    \setlength{\columnsep}{0.45cm}

    \setlength{\abovedisplayskip}{1.5pt}
    \setlength{\belowdisplayskip}{1.5pt}
    \setlength{\abovedisplayshortskip}{1pt}
    \setlength{\belowdisplayshortskip}{1pt}

    \setlength{\QuestionBefore}{0.4mm}
    \setlength{\QuestionAfter}{0.25mm}

    \setlength{\AnswerItemSep}{0pt}
    \setlength{\AnswerTopSep}{0.4mm}
    \setlength{\AnswerIndent}{0.4cm}
    \setlength{\InlineAnswerSep}{0.55cm}

    % Separación vertical entre opciones del formato 2x2
    \setlength{\GridAnswerSep}{2mm}
}

% =========================================================
% MODO ESTÁNDAR
% =========================================================

\newcommand{\ExamStandard}{
    \geometry{
        letterpaper,
        left=1.8cm,
        right=1.8cm,
        top=1.5cm,
        bottom=1.5cm
    }

    \setlength{\parindent}{0pt}
    \setlength{\parskip}{1.5pt}
    \setlength{\columnsep}{0.8cm}

    \setlength{\abovedisplayskip}{6pt}
    \setlength{\belowdisplayskip}{6pt}
    \setlength{\abovedisplayshortskip}{4pt}
    \setlength{\belowdisplayshortskip}{4pt}

    \setlength{\QuestionBefore}{2.5mm}
    \setlength{\QuestionAfter}{1.8mm}

    \setlength{\AnswerItemSep}{1.2mm}
    \setlength{\AnswerTopSep}{1.5mm}
    \setlength{\AnswerIndent}{0.55cm}
    \setlength{\InlineAnswerSep}{1cm}

    % Separación vertical entre opciones del formato 2x2
    \setlength{\GridAnswerSep}{5mm}
}

% =========================================================
% ENCABEZADO DEL EXAMEN
% =========================================================

% =========================================================
% ENCABEZADO DEL EXAMEN
% Estilo cercano al de las guías, pero compacto
% =========================================================

\newcommand{\ExamHeader}{
    \noindent
    \begin{minipage}[c]{0.65\textwidth}
        {\bfseries\large \ExamTitle}

        \vspace{0.5mm}

        {\bfseries \CourseName
            \hspace{1mm} \textbullet \hspace{1mm}
            \CourseCode}

        \vspace{0.5mm}

        {\small\color{ExamGray} \ExamTopics}

        \vspace{0.8mm}

        {\footnotesize
            Docente: \Instructor
            \hfill
            Periodo: \AcademicPeriod}

    \end{minipage}
    \hfill
    \begin{minipage}[c]{0.3\textwidth}
        \centering
        \MiniLogo
    \end{minipage}

    \vspace{1.5mm}

    {\color{ExamGreen}\hrule height 0.8pt}

    \vspace{1.3mm}

    \noindent
    {\footnotesize
        \textbf{Duración:} \ExamDuration
        \hfill
        \textbf{Valor:} \ExamValue
        \hfill
        \textbf{Preguntas:} \ExamQuestions
    }

    \vspace{2mm}
}

% =========================================================
% INFORMACIÓN DEL ESTUDIANTE
% =========================================================

\newcommand{\StudentInfo}{
    \noindent
    \textbf{Nombre:}
    \rule{7.4cm}{0.3pt}
    \hfill
    \textbf{Grupo:}
    \rule{3.0cm}{0.3pt}

    \vspace{2mm}

    \noindent
    \textbf{Documento:}
    \rule{7.0cm}{0.3pt}
    \hfill
    \textbf{Fecha:}
    \rule{3.0cm}{0.3pt}

    \vspace{2mm}
}

% =========================================================
% INSTRUCCIONES
% =========================================================

\newcommand{\ExamInstructions}[1]{
    \begin{adjustwidth}{0.2cm}{0.2cm}
        \small
        \color{ExamBlack}

        \textbf{Indicaciones.}
        #1
    \end{adjustwidth}

    \vspace{1.5mm}
}

% =========================================================
% CONTADOR DE PREGUNTAS
% =========================================================

\newcounter{examquestion}

\newenvironment{question}{
    \refstepcounter{examquestion}

    \par
    \vspace{\QuestionBefore}

    \noindent
    \textbf{\theexamquestion.}\,
    \ignorespaces
}{
    \par
    \vspace{\QuestionAfter}
}
% =========================================================
% RESPUESTAS VERTICALES
% =========================================================

\newenvironment{answers}{
    \begin{adjustwidth}{\AnswerIndent}{0cm}
        \begin{enumerate}[
            label=\textbf{\Alph*.},
            leftmargin=5mm,
            labelsep=1.5mm,
            itemsep=\AnswerItemSep,
            topsep=\AnswerTopSep,
            parsep=0pt,
            partopsep=0pt
            ]
        }{
        \end{enumerate}
    \end{adjustwidth}
}

% =========================================================
% RESPUESTAS EN UNA LÍNEA
% =========================================================

\newenvironment{inlineanswers}{
    \begin{adjustwidth}{\AnswerIndent}{0cm}
        \begin{enumerate*}[
            label=\textbf{\Alph*.},
            itemjoin=\hspace{\InlineAnswerSep},
            itemjoin*=\hspace{\InlineAnswerSep},
            itemsep=0pt
            ]
        }{
        \end{enumerate*}
    \end{adjustwidth}
}


% =========================================================
% RESPUESTAS EN CUADRÍCULA 2 x 2
% =========================================================

\newenvironment{gridanswers}{
    \begin{adjustwidth}{\AnswerIndent}{0cm}
        \begin{multicols}{2}
            \begin{enumerate}[
                label=\textbf{\Alph*.},
                leftmargin=5mm,
                labelsep=1.5mm,
                itemsep=\GridAnswerSep,
                topsep=\AnswerTopSep,
                parsep=0pt,
                partopsep=0pt
                ]
            }{
            \end{enumerate}
        \end{multicols}
    \end{adjustwidth}
}

% =========================================================
% COMANDO GENÉRICO PARA RESPUESTA
% =========================================================

\newcommand{\answer}{\item}

% =========================================================
% ESPACIO PARA PROCEDIMIENTO
% =========================================================

\newcommand{\ProcedureSpace}[1]{
    \vspace{\QuestionAfter}

    {\footnotesize\itshape
        Espacio para procedimiento:}

    \par
    \vspace{#1}
}

% =========================================================
% NOTA BREVE
% =========================================================

\newcommand{\ExamNote}[1]{
    \begin{adjustwidth}{0.3cm}{0.3cm}
        \small
        \color{ExamGray}
        \textit{#1}
    \end{adjustwidth}
}

% =========================================================
% SEPARADOR
% =========================================================

\newcommand{\ExamSeparator}{
    \vspace{1mm}
    {\color{ExamLine}\hrule}
    \vspace{1mm}
}

% =========================================================
% HOJA DE RESPUESTAS
% =========================================================
%
% Uso:
%
% \AnswerSheet{5}
% \AnswerSheet{10}
% \AnswerSheet{15}
% \AnswerSheet{20}
%
% Si no se necesita hoja de respuestas,
% simplemente no utilizar este comando.
% =========================================================

\newcommand{\AnswerSheet}[1]{

    \vspace{2mm}

    \begin{center}
        {\bfseries\color{ExamGreen}
            HOJA DE RESPUESTAS}
    \end{center}

    \vspace{-1mm}

    \begin{center}

        \renewcommand{\arraystretch}{1.25}
        \setlength{\tabcolsep}{4.5pt}

        % =================================================
        % 5 PREGUNTAS
        % =================================================

        \ifnum#1=5

        \begin{tabular}{c|cccc}
            \toprule
            \# & A & B & C & D \\
            \midrule

            1 & $\square$ & $\square$ & $\square$ & $\square$ \\
            2 & $\square$ & $\square$ & $\square$ & $\square$ \\
            3 & $\square$ & $\square$ & $\square$ & $\square$ \\
            4 & $\square$ & $\square$ & $\square$ & $\square$ \\
            5 & $\square$ & $\square$ & $\square$ & $\square$ \\

            \bottomrule
        \end{tabular}

        % =================================================
        % 10 PREGUNTAS
        % =================================================

        \else\ifnum#1=10

        \begin{tabular}{
                c|cccc
                @{\qquad\qquad}
                c|cccc
            }

            \toprule

            \# & A & B & C & D &
            \# & A & B & C & D \\

            \midrule

            1 & $\square$ & $\square$ & $\square$ & $\square$ &
            6 & $\square$ & $\square$ & $\square$ & $\square$ \\

            2 & $\square$ & $\square$ & $\square$ & $\square$ &
            7 & $\square$ & $\square$ & $\square$ & $\square$ \\

            3 & $\square$ & $\square$ & $\square$ & $\square$ &
            8 & $\square$ & $\square$ & $\square$ & $\square$ \\

            4 & $\square$ & $\square$ & $\square$ & $\square$ &
            9 & $\square$ & $\square$ & $\square$ & $\square$ \\

            5 & $\square$ & $\square$ & $\square$ & $\square$ &
            10 & $\square$ & $\square$ & $\square$ & $\square$ \\

            \bottomrule
        \end{tabular}

        % =================================================
        % 15 PREGUNTAS
        % =================================================

        \else\ifnum#1=15

        \begin{tabular}{
                c|cccc
                @{\qquad\qquad}
                c|cccc
            }

            \toprule

            \# & A & B & C & D &
            \# & A & B & C & D \\

            \midrule

            1 & $\square$ & $\square$ & $\square$ & $\square$ &
            9 & $\square$ & $\square$ & $\square$ & $\square$ \\

            2 & $\square$ & $\square$ & $\square$ & $\square$ &
            10 & $\square$ & $\square$ & $\square$ & $\square$ \\

            3 & $\square$ & $\square$ & $\square$ & $\square$ &
            11 & $\square$ & $\square$ & $\square$ & $\square$ \\

            4 & $\square$ & $\square$ & $\square$ & $\square$ &
            12 & $\square$ & $\square$ & $\square$ & $\square$ \\

            5 & $\square$ & $\square$ & $\square$ & $\square$ &
            13 & $\square$ & $\square$ & $\square$ & $\square$ \\

            6 & $\square$ & $\square$ & $\square$ & $\square$ &
            14 & $\square$ & $\square$ & $\square$ & $\square$ \\

            7 & $\square$ & $\square$ & $\square$ & $\square$ &
            15 & $\square$ & $\square$ & $\square$ & $\square$ \\

            8 & $\square$ & $\square$ & $\square$ & $\square$ &
            &           &           &           &           \\

            \bottomrule
        \end{tabular}

        % =================================================
        % 20 PREGUNTAS
        % =================================================

        \else\ifnum#1=20

        \begin{tabular}{
                c|cccc
                @{\qquad\qquad}
                c|cccc
            }

            \toprule

            \# & A & B & C & D &
            \# & A & B & C & D \\

            \midrule

            1  & $\square$ & $\square$ & $\square$ & $\square$ &
            11 & $\square$ & $\square$ & $\square$ & $\square$ \\

            2  & $\square$ & $\square$ & $\square$ & $\square$ &
            12 & $\square$ & $\square$ & $\square$ & $\square$ \\

            3  & $\square$ & $\square$ & $\square$ & $\square$ &
            13 & $\square$ & $\square$ & $\square$ & $\square$ \\

            4  & $\square$ & $\square$ & $\square$ & $\square$ &
            14 & $\square$ & $\square$ & $\square$ & $\square$ \\

            5  & $\square$ & $\square$ & $\square$ & $\square$ &
            15 & $\square$ & $\square$ & $\square$ & $\square$ \\

            6  & $\square$ & $\square$ & $\square$ & $\square$ &
            16 & $\square$ & $\square$ & $\square$ & $\square$ \\

            7  & $\square$ & $\square$ & $\square$ & $\square$ &
            17 & $\square$ & $\square$ & $\square$ & $\square$ \\

            8  & $\square$ & $\square$ & $\square$ & $\square$ &
            18 & $\square$ & $\square$ & $\square$ & $\square$ \\

            9  & $\square$ & $\square$ & $\square$ & $\square$ &
            19 & $\square$ & $\square$ & $\square$ & $\square$ \\

            10 & $\square$ & $\square$ & $\square$ & $\square$ &
            20 & $\square$ & $\square$ & $\square$ & $\square$ \\

            \bottomrule
        \end{tabular}

        % =================================================
        % VALOR NO CONFIGURADO
        % =================================================

        \else

        \fbox{
            \begin{minipage}{0.75\textwidth}
                \centering
                \small
                Hoja de respuestas no configurada para
                #1 preguntas.
            \end{minipage}
        }

        \fi
        \fi
        \fi
        \fi

    \end{center}
}

% =========================================================
% CALIFICACIÓN
% =========================================================
\newcommand{\ObservationLines}{%
    \par\vspace{3mm}
    \rule{\linewidth}{0.3pt}
    \par\vspace{3mm}
    \rule{\linewidth}{0.3pt}
}

\newcommand{\ExamScore}{
    \vspace{2mm}

    {\color{ExamLine}\hrule}

    \vspace{2mm}

    \noindent
    \textbf{Puntaje:}
    \rule{1.4cm}{0.3pt}
    /\ExamValue
    \hfill
    \textbf{Nota:}
    \rule{1.4cm}{0.3pt}
}

% =========================================================
% PÁGINA DE PROCEDIMIENTOS
% =========================================================

\newcommand{\ProcedurePage}{
    \newpage

    \begin{center}
        {\large\bfseries
            Procedimientos}
    \end{center}

    \vspace{1mm}

    \StudentInfo

    {\color{ExamGreen}\hrule height 0.6pt}

    \vspace{2mm}

    {\small
        Utilice este espacio para desarrollar los cálculos
        correspondientes a las preguntas que lo requieran.
    }

    \vfill

    \ExamScore
}

% =========================================================
% ESTILOS TIKZ PARA EXÁMENES DE ROS
% =========================================================

\tikzset{
    rosnode/.style={
        ellipse,
        draw=BrandBlue,
        thick,
        minimum width=2.5cm,
        minimum height=0.8cm,
        align=center
    },
    rostopic/.style={
        rectangle,
        draw=BrandYellow,
        thick,
        minimum width=2.8cm,
        minimum height=0.75cm,
        align=center
    },
    rosarrow/.style={
        -{Latex[length=2.5mm]},
        thick,
        BrandBlue
    }
}

% =========================================================
% FIN DEL TEMPLATE
% =========================================================

% =========================================================
% INFORMACIÓN DEL EXAMEN
% =========================================================

\renewcommand{\CourseName}{Robótica Industrial}
\renewcommand{\CourseCode}{ING 01339 / SOL-ING 220}
\renewcommand{\Institution}{Politécnico Colombiano Jaime Isaza Cadavid}
\renewcommand{\Faculty}{Facultad de Ingeniería}
\renewcommand{\Instructor}{Deimer Miranda Montoya, MSc.(c).}
\renewcommand{\ExamTitle}{Quiz \#1}
\renewcommand{\ExamTopics}{Fundamentación matemática, rotaciones 2D y marcos rotados}
\renewcommand{\AcademicPeriod}{2026-2}
\renewcommand{\ExamDuration}{60 minutos}
\renewcommand{\ExamValue}{100 puntos}
\renewcommand{\ExamQuestions}{10}
\renewcommand{\InstitutionLogo}{Template/LogoPoli.png}

% \ExamCompact
\ExamStandard

% =========================================================
\begin{document}

    \pagestyle{empty}

    % =========================================================
    % ENCABEZADO
    % =========================================================

    \ExamHeader
    \StudentInfo

    \ExamInstructions{
        Seleccione una única respuesta en cada pregunta. El valor de cada pregunta se
        indica en el enunciado. En las preguntas que involucren cálculo, el procedimiento
        podrá utilizarse como criterio de revisión. Use grados para los ángulos y conserve
        formas exactas con radicales cuando sea posible.
    }

    % =========================================================
    % PREGUNTAS
    % =========================================================

    \begin{multicols}{2}
        \raggedcolumns

        % ---------------------------------------------------------
        \begin{question}
            \textbf{[7 puntos]} Sea
            \[
            \mathbf{v}=\begin{bmatrix}3\\-4\\12\end{bmatrix}.
            \]
            ¿Cuál es su vector unitario asociado?

            \begin{inlineanswers}
                \answer{$\begin{bmatrix}3/13\\-4/13\\12/13\end{bmatrix}$}
                \answer{$\begin{bmatrix}3/12\\-4/12\\1\end{bmatrix}$}
                \answer{$\begin{bmatrix}3/17\\-4/17\\12/17\end{bmatrix}$}
                \answer{$\begin{bmatrix}-3/13\\4/13\\-12/13\end{bmatrix}$}
            \end{inlineanswers}
        \end{question}

        % RESPUESTA DOCENTE: A
        % Procedimiento:
        % ||v|| = sqrt(3^2+(-4)^2+12^2)
        %       = sqrt(9+16+144)
        %       = sqrt(169)
        %       = 13.
        % Por definición, v_hat = v/||v||.
        % Entonces:
        % v_hat = [3/13,-4/13,12/13]^T.
        % Verificación: ||v_hat|| = 1.

        % ---------------------------------------------------------
        \begin{question}
            \textbf{[8 puntos]} Sean
            \[
            \mathbf{a}=\begin{bmatrix}2\\1\\-1\end{bmatrix},
            \qquad
            \mathbf{b}=\begin{bmatrix}1\\-1\\1\end{bmatrix}.
            \]
            El ángulo entre ambos vectores es:

            \begin{inlineanswers}
                \answer{$0^\circ$}
                \answer{$45^\circ$}
                \answer{$90^\circ$}
                \answer{$180^\circ$}
            \end{inlineanswers}
        \end{question}

        % RESPUESTA DOCENTE: C
        % Procedimiento:
        % a·b = 2(1)+1(-1)+(-1)(1)
        %     = 2-1-1
        %     = 0.
        % Como a y b son no nulos y a·b=0, son ortogonales.
        % Por tanto, cos(theta)=0 y theta=90°.

        % ---------------------------------------------------------
        \begin{question}
            \textbf{[10 puntos]} Considere
            \[
            Q=
            \begin{bmatrix}
                \frac35 & \frac45\\[1mm]
                \frac45 & -\frac35
            \end{bmatrix}.
            \]
            ¿Cuál afirmación es correcta?

            \begin{answers}
                \answer{$Q$ es una rotación propia porque $Q^TQ=I$.}
                \answer{$Q$ es una reflexión: $Q^TQ=I$ y $\det(Q)=-1$.}
                \answer{$Q$ no es ortogonal porque $\det(Q)\neq 1$.}
                \answer{$Q$ es una rotación de $90^\circ$.}
            \end{answers}
        \end{question}

        % RESPUESTA DOCENTE: B
        % Procedimiento:
        % Las columnas son q1=[3/5,4/5]^T y q2=[4/5,-3/5]^T.
        % ||q1||=||q2||=1 y q1·q2=12/25-12/25=0.
        % Por tanto, Q^TQ=I: Q es ortogonal.
        % Sin embargo:
        % det(Q)=(3/5)(-3/5)-(4/5)(4/5)
        %       =-9/25-16/25
        %       =-1.
        % Una rotación propia requiere det(Q)=+1.
        % Por eso Q representa una reflexión, no una rotación propia.
                \columnbreak
        % ---------------------------------------------------------
        \begin{question}
            \textbf{[10 puntos]} El vector
            \[
            \mathbf{p}=\begin{bmatrix}2\\1\end{bmatrix}
            \]
            se rota $90^\circ$ en sentido antihorario. El resultado es:

            \begin{inlineanswers}
                \answer{$\begin{bmatrix}1\\2\end{bmatrix}$}
                \answer{$\begin{bmatrix}-1\\2\end{bmatrix}$}
                \answer{$\begin{bmatrix}1\\-2\end{bmatrix}$}
                \answer{$\begin{bmatrix}-2\\1\end{bmatrix}$}
            \end{inlineanswers}
        \end{question}

        % RESPUESTA DOCENTE: B
        % Procedimiento:
        % R(90°)=[[0,-1],[1,0]].
        % p' = R(90°)p
        %    = [[0,-1],[1,0]][2,1]^T
        %    = [-1,2]^T.
        % La norma se conserva: ||p||=sqrt(5)=||p'||.

        % ---------------------------------------------------------
        \begin{question}
            \textbf{[10 puntos]} La matriz
            \[
            R=
            \begin{bmatrix}
                -\frac{\sqrt3}{2} & -\frac12\\[1mm]
                \frac12 & -\frac{\sqrt3}{2}
            \end{bmatrix}
            \]
            corresponde a una rotación antihoraria de:

            \begin{inlineanswers}
                \answer{$30^\circ$}
                \answer{$120^\circ$}
                \answer{$150^\circ$}
                \answer{$210^\circ$}
            \end{inlineanswers}
        \end{question}

        % RESPUESTA DOCENTE: C
        % Procedimiento:
        % Para R(theta), la primera columna es [cos(theta),sin(theta)]^T.
        % Aquí:
        % cos(theta)=-sqrt(3)/2,
        % sin(theta)=1/2.
        % cos<0 y sin>0 => theta está en el segundo cuadrante.
        % El ángulo de referencia es 30°.
        % Entonces theta=180°-30°=150°.



        % ---------------------------------------------------------
        \begin{question}
            \textbf{[10 puntos]} Si
            \[
            R=R(30^\circ),
            \]
            ¿cuál expresión corresponde a $R^{-1}$?

            \begin{answers}
                \answer{$R^{-1}=R(30^\circ)$}
                \answer{$R^{-1}=R^T=R(-30^\circ)$}
                \answer{$R^{-1}=-R(30^\circ)$}
                \answer{$R^{-1}=I-R(30^\circ)$}
            \end{answers}
        \end{question}

        % RESPUESTA DOCENTE: B
        % Procedimiento:
        % Toda matriz de rotación propia cumple R^TR=I.
        % Multiplicando por R^{-1}, se obtiene R^{-1}=R^T.
        % En 2D, transponer R(theta) cambia el signo de los términos seno:
        % R(theta)^T = R(-theta).
        % Por tanto:
        % R(30°)^-1 = R(30°)^T = R(-30°).
                \columnbreak
        % ---------------------------------------------------------
        \begin{question}
            \textbf{[10 puntos]} Un vector inicialmente coincide con el eje $x$ positivo.
            Primero se rota $40^\circ$ en sentido antihorario y luego $100^\circ$ en
            sentido horario. La rotación equivalente y el vector final son:

            \begin{center}
                \begin{tikzpicture}[scale=0.78,>=Latex,every node/.style={font=\scriptsize}]
                    \draw[->,thick] (-0.2,0) -- (2.5,0) node[right] {$x$};
                    \draw[->,thick] (0,-2.1) -- (0,1.8) node[above] {$y$};

                    \draw[->,very thick,BrandBlue]
                    (0,0) -- (1.7,0)
                    node[above right] {$\mathbf{e}_x$};

                    \draw[->,very thick,BrandGreen]
                    (0,0) -- ({1.7*cos(-60)},{1.7*sin(-60)})
                    node[below right] {$\mathbf{p}_f$};

                    \draw[->,gray,thick]
                    (0.75,0) arc[start angle=0,end angle=-60,radius=0.75];
                    \node[gray] at (1.2,-0.38) {$-60^\circ$};
                \end{tikzpicture}
            \end{center}

            \begin{answers}
                \answer{$R(60^\circ)$ y $\begin{bmatrix}\frac12\\ \frac{\sqrt3}{2}\end{bmatrix}$}
                \answer{$R(-60^\circ)$ y $\begin{bmatrix}\frac12\\ -\frac{\sqrt3}{2}\end{bmatrix}$}
                \answer{$R(-140^\circ)$ y $\begin{bmatrix}-\frac{\sqrt3}{2}\\ \frac12\end{bmatrix}$}
                \answer{$R(140^\circ)$ y $\begin{bmatrix}-\frac12\\ \frac{\sqrt3}{2}\end{bmatrix}$}
            \end{answers}
        \end{question}

        % RESPUESTA DOCENTE: B
        % Procedimiento:
        % Primer giro: +40°.
        % Segundo giro: -100°.
        % En 2D las rotaciones sucesivas suman sus ángulos:
        % theta_eq = 40° + (-100°) = -60°.
        % Como el vector inicial es e_x=[1,0]^T:
        % p_f=R(-60°)e_x
        %    =[cos(-60°), sin(-60°)]^T
        %    =[1/2,-sqrt(3)/2]^T.

        % ---------------------------------------------------------
        \begin{question}
            \textbf{[12 puntos]} Los marcos $\{A\}$ y $\{B\}$ comparten origen. El marco
            $\{B\}$ está rotado $30^\circ$ en sentido antihorario respecto a $\{A\}$ y
            \[
            {}^B\mathbf{p}=\begin{bmatrix}2\\1\end{bmatrix}.
            \]
            ¿Cuáles son las coordenadas del mismo punto expresadas en $\{A\}$?

            \begin{answers}
                \answer{$\begin{bmatrix}\sqrt3-\frac12\\[1mm]1+\frac{\sqrt3}{2}\end{bmatrix}$}
                \answer{$\begin{bmatrix}\sqrt3+\frac12\\[1mm]-1+\frac{\sqrt3}{2}\end{bmatrix}$}
                \answer{$\begin{bmatrix}2\sqrt3-1\\[1mm]2+\sqrt3\end{bmatrix}$}
                \answer{$\begin{bmatrix}2\\1\end{bmatrix}$}
            \end{answers}
        \end{question}

        % RESPUESTA DOCENTE: A
        % Procedimiento:
        % La notación A_R_B describe la orientación de B vista desde A.
        % Como B está rotado +30° respecto a A:
        % A_R_B=R(30°).
        % Para pasar coordenadas de B hacia A:
        % p_A = A_R_B p_B.
        % Entonces:
        % p_A = [[cos30,-sin30],[sin30,cos30]][2,1]^T
        %     = [2cos30-sin30, 2sin30+cos30]^T
        %     = [sqrt(3)-1/2, 1+sqrt(3)/2]^T.
        % La opción C es exactamente el doble del resultado y es un distractor.

        % ---------------------------------------------------------
        \begin{question}
            \textbf{[12 puntos]} El marco $\{B\}$ está rotado $45^\circ$ en sentido
            antihorario respecto a $\{A\}$. Si
            \[
            {}^A\mathbf{q}=\begin{bmatrix}0\\2\end{bmatrix},
            \]
            ¿cuáles son las coordenadas de ese punto expresadas en $\{B\}$?

            \begin{inlineanswers}
                \answer{$\begin{bmatrix}\sqrt2\\\sqrt2\end{bmatrix}$}
                \answer{$\begin{bmatrix}-\sqrt2\\\sqrt2\end{bmatrix}$}
                \answer{$\begin{bmatrix}\sqrt2\\-\sqrt2\end{bmatrix}$}
                \answer{$\begin{bmatrix}0\\2\end{bmatrix}$}
            \end{inlineanswers}
        \end{question}

        % RESPUESTA DOCENTE: A
        % Procedimiento:
        % Se conoce A_R_B=R(45°).
        % Para pasar de coordenadas en A hacia coordenadas en B se utiliza la inversa:
        % p_B = (A_R_B)^-1 p_A.
        % Como la matriz es de rotación:
        % (A_R_B)^-1=(A_R_B)^T=R(-45°).
        % Entonces:
        % q_B = R(-45°)[0,2]^T
        %     = [sqrt(2),sqrt(2)]^T.

        % ---------------------------------------------------------
        \begin{question}
            \textbf{[11 puntos]} Observe los marcos $\{A\}$ y $\{B\}$:

            \begin{center}
                \begin{tikzpicture}[scale=0.78,>=Latex,every node/.style={font=\scriptsize}]
                    \coordinate (O) at (0,0);

                    \draw[->,very thick,BrandBlue]
                    (O) -- (2.5,0) node[right] {$x_A$};
                    \draw[->,very thick,BrandBlue]
                    (O) -- (0,2.5) node[above] {$y_A$};

                    \draw[->,very thick,BrandGreen]
                    (O) -- (0,2.0) node[right] {$x_B$};
                    \draw[->,very thick,BrandGreen]
                    (O) -- (-2.0,0) node[above] {$y_B$};

                    \draw[->,gray,thick]
                    (0.65,0) arc[start angle=0,end angle=90,radius=0.65];
                    \node[gray] at (1.0,0.5) {$90^\circ$};
                \end{tikzpicture}
            \end{center}

            Para esta configuración,
            \[
            {}^A R_B=
            \begin{bmatrix}
                0 & -1\\
                1 & 0
            \end{bmatrix}.
            \]
            ¿Cuál interpretación es correcta?

            \begin{answers}
                \answer{La primera columna indica que $x_B$ coincide con $-y_A$.}
                \answer{La segunda columna indica que $y_B$ coincide con $+x_A$.}
                \answer{$x_B$ coincide con $+y_A$, $y_B$ con $-x_A$ y para pasar de $A$ a $B$ se usa $({}^A R_B)^T$.}
                \answer{El punto físico debe rotarse $90^\circ$ cada vez que cambia de marco.}
            \end{answers}
        \end{question}

        % RESPUESTA DOCENTE: C
        % Procedimiento:
        % Las columnas de A_R_B son los ejes de B expresados desde A.
        % Primera columna [0,1]^T:
        % x_B coincide con +y_A.
        % Segunda columna [-1,0]^T:
        % y_B coincide con -x_A.
        % Además, para convertir coordenadas de A hacia B:
        % p_B = (A_R_B)^-1 p_A.
        % Como A_R_B es una rotación propia:
        % (A_R_B)^-1 = (A_R_B)^T.
        % El punto físico no gira por cambiar de sistema de referencia; sólo cambia su
        % representación numérica.

    \end{multicols}

    % =========================================================
    % HOJA DE RESPUESTAS Y CALIFICACIÓN
    % =========================================================

    \AnswerSheet{10}

    \vspace{2mm}
    {\color{ExamGreen}\hrule height 0.6pt}
    \vspace{2mm}

    \ExamScore

\end{document}
